% GENERATED FILE. Edit canonical lessons, then run npm run generate:books.
% canonical-source-hash: fnv1a64:0a4901d8f9c6ba05
% canonical-lessons: ML-C541-samikkuka, ML-C541-vasalakuka, ML-C541-pakaruka, ML-C541-upavasikkuka, ML-C541-dahikkuka

\chapter{Getting Worse, Getting Better}
\label{ch:ml-getting-worse-getting-better}
\hlchaptermodality{541}

\begin{chapteropening}
\textbf{By the end of this chapter:} \emph{I can say to subside, to get worse, to spread by contact, to go without food and to be digested.}

Complete the last lesson of chapter 541: I can say to subside, to get worse, to spread by contact, to go without food and to be digested.
\end{chapteropening}

\section[\texorpdfstring{\ml{ശമിക്കുക} (śamikkuka)}{ശമിക്കുക (śamikkuka)}]{\ml{ശമിക്കുക} (śamikkuka) — to subside; to ease}

\label{lesson:ML-C541-samikkuka}



\begin{quote}
Before the new one: say the Malayalam for to go numb, then the Malayalam for to feel sore.
\end{quote}

\subsection*{You'll want to know: \ml{ശമിക്കുക}}
\textbf{\ml{ശമിക്കുക}} — \emph{śamikkuka} — \textquotedblleft{}to subside; to ease\textquotedblright{}.

\begin{quote}
\ml{മരുന്ന്} \ml{കഴിച്ചപ്പോൾ} \ml{വേദന} \ml{ശമിച്ചു}.
\end{quote}

\begin{quote}
\emph{marunnŭ kaḻiccappōḷ vēdana śamiccu.}
\end{quote}

\begin{quote}
The pain eased when I took the medicine.
\end{quote}

Past: \ml{ശമിച്ചു}, \emph{śamiccu}.

Formal; also said of rain, anger and fire.

\begin{cousinweb}
\textbf{\ml{ശമിക്കുക}} (\emph{śamikkuka}) is built from Sanskrit \emph{śama}, \textquotedblleft{}tranquillity, calm\textquotedblright{} (Monier-Williams), with the Malayalam verb ending -\ml{ിക്കുക}.
\end{cousinweb}

\begin{grammarlens}[title={What you've built}]
One more word for describing what you lived and felt.
\end{grammarlens}

\subsection*{Your turn}
\begin{itemize}
\raggedright
  \item \emph{Say it:} \emph{śamikkuka}
  \item \emph{Say it:} \emph{śamikkuka}, once more
  \item \emph{Say it:} \emph{marunnŭ kaḻiccappōḷ vēdana śamiccu.} — The pain eased when I took the medicine.
  \item \emph{Recall:} say the Malayalam for to go numb, then the Malayalam for to feel sore, then say \emph{śamikkuka} again
\end{itemize}

\begin{tcolorbox}[breakable,colback=teal!4,colframe=teal!35!black,title={Before you move on}]
What does \ml{ശമിക്കുക} mean? (\textquotedblleft{}To subside; to ease\textquotedblright{}.) Say it once more.
\end{tcolorbox}
\section[\texorpdfstring{\ml{വഷളാകുക} (vaṣaḷākuka)}{വഷളാകുക (vaṣaḷākuka)}]{\ml{വഷളാകുക} (vaṣaḷākuka) — to get worse; to worsen}

\label{lesson:ML-C541-vasalakuka}



\begin{quote}
Before the new one: say the Malayalam for to feel sore, then the Malayalam for to subside.
\end{quote}

\subsection*{You'll want to know: \ml{വഷളാകുക}}
\textbf{\ml{വഷളാകുക}} — \emph{vaṣaḷākuka} — \textquotedblleft{}to get worse; to worsen\textquotedblright{}.

\begin{quote}
\ml{ചികിത്സ} \ml{വൈകിയതുകൊണ്ട്} \ml{രോഗം} \ml{വഷളായി}.
\end{quote}

\begin{quote}
\emph{cikitsa vaikiyatukoṇṭŭ rōgaṁ vaṣaḷāyi.}
\end{quote}

\begin{quote}
The illness got worse because treatment was late.
\end{quote}

Past: \ml{വഷളായി}, \emph{vaṣaḷāyi}.

\ml{വഷളാക്കുക} is to make worse, or to spoil a child.

\begin{cousinweb}
\textbf{\ml{വഷളാകുക}} (\emph{vaṣaḷākuka}) is an everyday word; no source we can open explains it, so its origin is not settled.
\end{cousinweb}

\begin{grammarlens}[title={What you've built}]
One more word for describing what you lived and felt.
\end{grammarlens}

\subsection*{Your turn}
\begin{itemize}
\raggedright
  \item \emph{Say it:} \emph{vaṣaḷākuka}
  \item \emph{Say it:} \emph{vaṣaḷākuka}, once more
  \item \emph{Say it:} \emph{cikitsa vaikiyatukoṇṭŭ rōgaṁ vaṣaḷāyi.} — The illness got worse because treatment was late.
  \item \emph{Recall:} say the Malayalam for to feel sore, then the Malayalam for to subside, then say \emph{vaṣaḷākuka} again
\end{itemize}

\begin{tcolorbox}[breakable,colback=teal!4,colframe=teal!35!black,title={Before you move on}]
What does \ml{വഷളാകുക} mean? (\textquotedblleft{}To get worse; to worsen\textquotedblright{}.) Say it once more.
\end{tcolorbox}
\section[\texorpdfstring{\ml{പകരുക} (pakaruka)}{പകരുക (pakaruka)}]{\ml{പകരുക} (pakaruka) — to spread by contact; to be passed on}

\label{lesson:ML-C541-pakaruka}



\begin{quote}
Before the new one: say the Malayalam for to subside, then the Malayalam for to get worse.
\end{quote}

\subsection*{You'll want to know: \ml{പകരുക}}
\textbf{\ml{പകരുക}} — \emph{pakaruka} — \textquotedblleft{}to spread by contact; to be passed on\textquotedblright{}.

\begin{quote}
\ml{ജലദോഷം} \ml{ഒരാളിൽ} \ml{നിന്ന്} \ml{മറ്റൊരാൾക്ക്} \ml{പകരും}.
\end{quote}

\begin{quote}
\emph{jaladōṣaṁ orāḷil ninnŭ maṟṟorāḷkkŭ pakaruṁ.}
\end{quote}

\begin{quote}
A cold spreads from one person to another.
\end{quote}

Past: \ml{പകർന്നു}, \emph{pakarnnu}.

It also means to pour: \ml{അമ്മ} \ml{ചായ} \ml{ഗ്ലാസിലേക്ക്} \ml{പകർന്നു}.

\begin{cousinweb}
\textbf{\ml{പകരുക}} (\emph{pakaruka}) is native Dravidian: DEDR 3803 gives it as to change place and to be exchanged (DEDR 3803). An illness that is caught changes place from one person to the next.
\end{cousinweb}

\begin{grammarlens}[title={What you've built}]
One more word for describing what you lived and felt.
\end{grammarlens}

\subsection*{Your turn}
\begin{itemize}
\raggedright
  \item \emph{Say it:} \emph{pakaruka}
  \item \emph{Say it:} \emph{pakaruka}, once more
  \item \emph{Say it:} \emph{jaladōṣaṁ orāḷil ninnŭ maṟṟorāḷkkŭ pakaruṁ.} — A cold spreads from one person to another.
  \item \emph{Recall:} say the Malayalam for to subside, then the Malayalam for to get worse, then say \emph{pakaruka} again
\end{itemize}

\begin{tcolorbox}[breakable,colback=teal!4,colframe=teal!35!black,title={Before you move on}]
What does \ml{പകരുക} mean? (\textquotedblleft{}To spread by contact; to be passed on\textquotedblright{}.) Say it once more.
\end{tcolorbox}
\section[\texorpdfstring{\ml{ഉപവസിക്കുക} (upavasikkuka)}{ഉപവസിക്കുക (upavasikkuka)}]{\ml{ഉപവസിക്കുക} (upavasikkuka) — to go without food; to fast}

\label{lesson:ML-C541-upavasikkuka}



\begin{quote}
Before the new one: say the Malayalam for to get worse, then the Malayalam for to spread by contact.
\end{quote}

\subsection*{You'll want to know: \ml{ഉപവസിക്കുക}}
\textbf{\ml{ഉപവസിക്കുക}} — \emph{upavasikkuka} — \textquotedblleft{}to go without food; to fast\textquotedblright{}.

\begin{quote}
\ml{രക്തപരിശോധനയ്ക്ക്} \ml{മുൻപ്} \ml{ഞാൻ} \ml{രാത്രി} \ml{ഉപവസിച്ചു}.
\end{quote}

\begin{quote}
\emph{raktapariśōdhanaykkŭ munpŭ ñān rātri upavasiccu.}
\end{quote}

\begin{quote}
I fasted overnight before the blood test.
\end{quote}

Past: \ml{ഉപവസിച്ചു}, \emph{upavasiccu}.

The noun is \ml{ഉപവാസം}, a fast.

\begin{cousinweb}
\textbf{\ml{ഉപവസിക്കുക}} (\emph{upavasikkuka}) is built from Sanskrit \emph{upavāsa}, \textquotedblleft{}a fast, fasting\textquotedblright{} (CDIAL 2244), with the Malayalam verb ending -\ml{ിക്കുക}.
\end{cousinweb}

\begin{grammarlens}[title={What you've built}]
One more word for describing what you lived and felt.
\end{grammarlens}

\subsection*{Your turn}
\begin{itemize}
\raggedright
  \item \emph{Say it:} \emph{upavasikkuka}
  \item \emph{Say it:} \emph{upavasikkuka}, once more
  \item \emph{Say it:} \emph{raktapariśōdhanaykkŭ munpŭ ñān rātri upavasiccu.} — I fasted overnight before the blood test.
  \item \emph{Recall:} say the Malayalam for to get worse, then the Malayalam for to spread by contact, then say \emph{upavasikkuka} again
\end{itemize}

\begin{tcolorbox}[breakable,colback=teal!4,colframe=teal!35!black,title={Before you move on}]
What does \ml{ഉപവസിക്കുക} mean? (\textquotedblleft{}To go without food; to fast\textquotedblright{}.) Say it once more.
\end{tcolorbox}
\section[\texorpdfstring{\ml{ദഹിക്കുക} (dahikkuka)}{ദഹിക്കുക (dahikkuka)}]{\ml{ദഹിക്കുക} (dahikkuka) — to be digested; to digest}

\label{lesson:ML-C541-dahikkuka}



\begin{quote}
Before the new one: say the Malayalam for to spread by contact, then the Malayalam for to go without food.
\end{quote}

\subsection*{You'll want to know: \ml{ദഹിക്കുക}}
\textbf{\ml{ദഹിക്കുക}} — \emph{dahikkuka} — \textquotedblleft{}to be digested; to digest\textquotedblright{}.

\begin{quote}
\ml{ഇറച്ചി} \ml{എളുപ്പം} \ml{ദഹിക്കില്ല}.
\end{quote}

\begin{quote}
\emph{iṟacci eḷuppaṁ dahikkilla.}
\end{quote}

\begin{quote}
Meat does not digest easily.
\end{quote}

Past: \ml{ദഹിച്ചു}, \emph{dahiccu}.

The noun is \ml{ദഹനം}, digestion; \ml{ദഹനക്കേട്} is indigestion.

\begin{cousinweb}
\textbf{\ml{ദഹിക്കുക}} (\emph{dahikkuka}) is built from Sanskrit \emph{dahati}, \textquotedblleft{}burns\textquotedblright{} (CDIAL 6245), with the Malayalam verb ending -\ml{ിക്കുക}: in Malayalam the burning up of food became digesting.
\end{cousinweb}

\begin{grammarlens}[title={What you've built}]
One more word for describing what you lived and felt.
\end{grammarlens}

\subsection*{Your turn}
\begin{itemize}
\raggedright
  \item \emph{Say it:} \emph{dahikkuka}
  \item \emph{Say it:} \emph{dahikkuka}, once more
  \item \emph{Say it:} \emph{iṟacci eḷuppaṁ dahikkilla.} — Meat does not digest easily.
  \item \emph{Recall:} say the Malayalam for to spread by contact, then the Malayalam for to go without food, then say \emph{dahikkuka} again
\end{itemize}

\begin{tcolorbox}[breakable,colback=teal!4,colframe=teal!35!black,title={Before you move on}]
What does \ml{ദഹിക്കുക} mean? (\textquotedblleft{}To be digested; to digest\textquotedblright{}.) Say it once more.
\end{tcolorbox}
